\section{Related Work}
Cross-lingual NER is a crowded area, especially for low-resource targets. English-to-Danish transfer has already been shown to work as a practical zero-shot baseline, improving quickly once even modest Danish supervision is added \citep{plank2020danish}. Our results broadly corroborate that picture while adding a direct Norwegian comparison. More broadly, work such as \citet{fu2023heterogeneous} finds that structural similarity often predicts transfer success — a claim we pressure-test here with a controlled setup that holds architecture, data format, and evaluation protocol constant.

Two Danish-specific resources anchor the fully supervised baseline literature. DaNE \citep{hvingelby2020dane} provides the manually annotated corpus that underlies most subsequent Danish NER work, and DaCy \citep{enevoldsen2021dacy} demonstrates strong Danish-supervised pipelines built on the same annotation conventions. Our Danish monolingual baselines are most directly comparable to that line of work.

Multilingual encoders make zero-shot transfer practically viable. BERT established the transformer pretraining paradigm \citep{devlin2019bert}, and XLM-R extended it with a substantially larger and more linguistically diverse corpus \citep{conneau2020xlmr}, often yielding better cross-lingual generalization. Whether that edge persists when source and target are as close as Norwegian and Danish is an open question that motivates our design.

Typological distance as a predictor of transfer performance has been explored via resources such as URIEL+ \citep{khan2025urielplus} and DistALS \citep{goot2025distals}; recent studies argue that surface similarity can explain a substantial portion of the zero-shot gap \citep{sohn2024zeroshot}. Our contribution is to compute WALS-feature distances explicitly and tie those measurements directly to observed F1 differences, using the Universal NER benchmark \citep{mayhew2024universal} across all three languages to keep annotation-format differences out of the picture.
